% Exercice: Serie S
% Sujet: Géométrie
% Généré automatiquement

\begin{exercice}[Géométrie]
\label{ex:60}

1- Déterminer l'arrondi d'ordre 2 de la moyenne des dis$tances hebdomadaires$\\
0,7^\{5\} pt\\
parcourues par ces élèves\\
Déterminer la classe modale de cette série statistique.\\
0,5 pt\\
Dresser le tableau des effectifs cumules croissants\\
0,7^\{5\}pt\\
Déterminer par interpolation linéaire\\
la médiane de cette série statistique\\
nombre délèves qui parcourent moins de 5 km par semaine\\
Déterminer\\
0,5 pt\\
le professeur titulaire de\\
mieux\\
En vue de\\
cette classe voudrait constituer des groupes détude de cinq élèves,\\
Pour chacune des deux questions suivantes, quatre réponses sont proposées parmi\\
suivi\\
lesquelles une seule est juste. Ecrire\\
numéro de la question\\
lettre\\
sur votre feuille de composition\\
Aucune justification\\
juste\\
correspondant à la réponse\\
n'est demandée\\
Le nombre de groupes possibles qu'il peut former est\\
0 Spt\\
9^\{0\}-\\
\item b) A^\{3\}o
Le nombre de groupes qu'il peut former contenant au moins deux filles et au moins\\
deux garçons est\\
C^\{3\}o x C^\{3\}o\\
C^\{3\}o * C^\{3\}o\\
C^\{3\}o\\
C^\{3\}o x C^\{3\}o\\
Abo X A^\{3\}o\\
\item d) C^\{3\}o
Ipt\\

\end{exercice}

% --- Fin Serie S ---
